\uuid{OGvR}
\exo7id{3909}
\titre{exo7 3909}
\auteur{quercia}
\organisation{exo7}
\datecreate{2010-03-11}
\isIndication{false}
\isCorrection{false}
\chapitre{Continuité, limite et étude de fonctions réelles}
\sousChapitre{Etude de fonctions}
\cpge{mpsi-mp2i.calcul-algebrique-trigonometrie.d, pcsi.calcul-algebrique-trigonometrie.d, ptsi.calcul-algebrique-trigonometrie.d, mpsi-mp2i.nombres-reels-suites.i, pcsi.nombres-reels-suites.g, ptsi.nombres-reels-suites.g, mpsi-mp2i.limites-continuite.b, pcsi.limites-continuite.b, ptsi.limites-continuite.b, bcpst1.trigonometrie, bcpst1.suites-reelles-usuelles, bcpst1.limites-continuite.c}
\module{Analyse}
\niveau{L1}
\difficulte{}

\contenu{
\question{
Montrer qu'il existe $\theta \in {]0, \frac \pi2 [}$ tel que pour tout
$\alpha \in \R$, les nombres
$\sin^2(\theta-\alpha), \sin^2\theta, \sin^2(\theta+\alpha)$
soient en progression arithmétique.
}
}
